\documentclass[conference]{IEEEtran}
\IEEEoverridecommandlockouts

% Standard IEEE Packages
\usepackage{cite}
\usepackage{amsmath,amssymb,amsfonts}
\usepackage{algorithmic}
\usepackage{algorithm}
\usepackage{graphicx}
\usepackage{textcomp}
\usepackage{xcolor}
\usepackage{booktabs}
\usepackage{url}
\usepackage{balance}
\usepackage{multirow}
\usepackage{microtype}

\def\BibTeX{{\rm B\kern-.05em{\sc i\kern-.025em b}\kern-.08em
    T\kern-.1667em\lower.7ex\hbox{E}\kern-.125emX}}

\setlength{\abovedisplayskip}{4pt plus 1pt minus 1pt}
\setlength{\belowdisplayskip}{4pt plus 1pt minus 1pt}
\setlength{\textfloatsep}{6pt plus 1pt minus 1pt}
\setlength{\floatsep}{6pt plus 1pt minus 1pt}
\setlength{\intextsep}{6pt plus 1pt minus 1pt}

\begin{document}

\title{FINDORA AI: An Intelligent Campus Lost and Found System Using Multimodal Matching and Smart Verification\\
\large \textbf{Track:} Smart Lost and Found Management System \quad|\quad \textbf{Author:} Tharunkumar K (@Tharun4743)
}

\author{
    \IEEEauthorblockN{Tharunkumar K}
    \IEEEauthorblockA{\textit{Dept. of Information Technology} \\
    \textit{V.S.B. Engineering College}\\
    Karur, Tamil Nadu, India \\
    tharunkumark42007@gmail.com}
}

\maketitle

\begin{abstract}
Managing lost and found items in educational institutions is a persistent challenge. Most college campuses still rely on manual notebook registers and physical notice boards. This traditional approach leads to low recovery rates (under 18\%), delayed item returns, and privacy risks where false claimants can easily guess item details. In this paper, we present \textbf{FINDORA AI}, an automated smart lost-and-found system developed for college campuses. Our system combines five key components: (1) \textit{Live Camera Evidence Capture} with GPS coordinates and real-time timestamp watermarking; (2) \textit{Multimodal Hybrid Matching} that evaluates text similarity (TF-IDF cosine similarity), visual color distance (normalized RGB metric), categorical alignment, spatial proximity (Haversine distance), and time decay; (3) \textit{Blind Challenge Verification} where claimants must answer dynamic questions about hidden item attributes before claiming; (4) \textit{Anti-Fraud Velocity Limiting} that flags rapid or suspicious claiming attempts; and (5) \textit{Secure One-Time Handover Codes} sent to verified student emails to ensure safe physical item collection. During testing at V.S.B. Engineering College (VSBEC), FINDORA AI achieved a \textbf{94.2\% Top-3 retrieval precision}, reduced average item recovery time from \textbf{72.4 hours to 4.2 hours} (a 17.2$\times$ speedup), and virtually eliminated fraudulent claims.
\end{abstract}

\begin{IEEEkeywords}
Campus Lost and Found, Multimodal Matching, TF-IDF Cosine Similarity, Image Watermarking, Blind Verification, Telegram Bot, PostgreSQL, Cloudinary.
\end{IEEEkeywords}

\section{Introduction}
In modern colleges and universities with thousands of students and faculty members, personal belongings like ID cards, laptops, mobile phones, headphones, keys, and notebooks are frequently misplaced in classrooms, libraries, laboratories, cafeterias, and buses.

Currently, most campuses handle lost property using traditional manual procedures:
\begin{itemize}
    \item \textbf{Paper Logbooks:} Finders and security desks write item descriptions by hand in physical registers, which are difficult to search and maintain.
    \item \textbf{Vague Descriptions:} When owners report an item, they often use different words than the finder (e.g., ``blue water bottle'' vs. ``navy sports flask''). Simple keyword searches in databases fail to match these items.
    \item \textbf{Privacy Exposure and False Claims:} Displaying lost item photos and descriptions on open notice boards or group chats allows unauthorized individuals to claim items that do not belong to them.
    \item \textbf{Unorganized Handover:} There is no digital proof or authentication mechanism during the physical handover process.
\end{itemize}

As a result of these limitations, statistical studies show that less than 18.2\% of lost items on college campuses are successfully returned to their rightful owners:
\begin{equation}
    \text{Recovery Rate}_{\text{manual}} = \frac{N_{\text{recovered}}}{N_{\text{reported}}} \le 18.2\%
\end{equation}

To solve this problem, Team Techsquad (Team 22) developed \textbf{FINDORA AI}. The goal of FINDORA AI is to provide a complete, automated, and secure digital platform that connects item losers and finders, accurately suggests matching items using multimodal algorithms, verifies true ownership through interactive blind questions, and coordinates safe pickup via automated email codes and Telegram alerts.

\section{System Architecture}
FINDORA AI is designed as a responsive, modular web application with cloud storage and asynchronous notification services.

\begin{figure}[htbp]
\centering
\small
\setlength{\tabcolsep}{4pt}
\begin{tabular}{|p{0.92\columnwidth}|}
\hline
\textbf{FINDORA AI System Architecture} \\
\hline
\textbf{1. User Interface (Client):} React 19 Single Page Application, TypeScript, Tailwind CSS, HTML5 Canvas live camera with GPS watermarking. \\
\textbf{2. Backend Server:} Node.js with Express REST APIs, JWT authentication, and Brevo SMTP email service. \\
\textbf{3. Database Storage:} Supabase PostgreSQL database for structured item records and user accounts. \\
\textbf{4. Media Storage:} Cloudinary CDN for verified evidence image hosting. \\
\textbf{5. Community Bot:} Telegram Bot API (\texttt{@findoravsb\_bot}) for real-time broadcast and instant queries. \\
\hline
\end{tabular}
\caption{Overall Architecture of the FINDORA AI Platform.}
\label{fig:architecture}
\end{figure}

\subsection{Client-Side Live Camera and Watermarking}
To ensure that reported items are authentic and physically present, FINDORA AI includes a built-in camera capture tool built with WebRTC. When an item is photographed:
\begin{enumerate}
    \item The device retrieves current GPS coordinates $(\text{Lat}, \text{Lon})$ and UTC timestamp.
    \item An HTML5 Canvas automatically stamps an indelible watermark containing the date, time, campus zone, and coordinates across the bottom of the photo.
    \item The watermarked image is uploaded directly to Cloudinary storage, generating a secure URL.
\end{enumerate}
This mechanism prevents users from submitting arbitrary internet images or fake photos.

\subsection{Database and Backend Gateway}
The backend is built with Node.js and Express. It connects to a PostgreSQL database hosted on Supabase, which stores item details, user profiles, blind challenge questions, and claim statuses with foreign key constraints. Sensitive private attributes (such as serial numbers or custom engravings) are kept confidential and are never displayed publicly.

\section{Multimodal Matching Algorithm}
When a user submits a lost or found report, FINDORA AI scans the existing database records of opposite type (i.e., comparing a ``lost'' report against all active ``found'' items) and calculates a comprehensive matching score $\Phi \in [0, 1]$.

\subsection{Mathematical Model}
The composite matching score $\Phi(I_t, I_j)$ between target item $I_t$ and candidate item $I_j$ is calculated as a weighted sum of five components:
\begin{equation}
    \Phi(I_t, I_j) = w_v S_{\text{vis}} + w_s S_{\text{sem}} + w_c S_{\text{cat}} + w_l S_{\text{spa}} + w_t S_{\text{tem}}
\end{equation}
where the weights satisfy $w_v + w_s + w_c + w_l + w_t = 1.0$.

Based on experimental tuning with campus test cases, the weights are assigned as:
\begin{itemize}
    \item \textbf{Visual Color Similarity ($w_v = 0.35$):} Compares primary item colors in normalized RGB space:
    \begin{equation}
        S_{\text{vis}} = 1.0 - \frac{\sqrt{(R_1-R_2)^2 + (G_1-G_2)^2 + (B_1-B_2)^2}}{255\sqrt{3}}
    \end{equation}
    \item \textbf{Text Semantic Similarity ($w_s = 0.30$):} Uses TF-IDF vectorization and cosine similarity over item titles and descriptions:
    \begin{equation}
        S_{\text{sem}} = \frac{\vec{v}_t \cdot \vec{v}_j}{\|\vec{v}_t\|_2 \|\vec{v}_j\|_2}
    \end{equation}
    \item \textbf{Category Concordance ($w_c = 0.15$):} Checks if the selected categories match:
    \begin{equation}
        S_{\text{cat}} = \begin{cases} 1.0, & \text{Category}(I_t) = \text{Category}(I_j) \\ 0.0, & \text{Otherwise} \end{cases}
    \end{equation}
    \item \textbf{Spatial Proximity ($w_l = 0.10$):} Calculates distance $\Delta d$ (in meters) using the Haversine formula and applies an exponential decay:
    \begin{equation}
        S_{\text{spa}} = \exp\left(-\frac{\Delta d}{\sigma_d}\right), \quad \sigma_d = 250\text{ m}
    \end{equation}
    \item \textbf{Temporal Decay ($w_t = 0.10$):} Measures time difference $\Delta t$ (in hours) between reported times:
    \begin{equation}
        S_{\text{tem}} = \exp\left(-\lambda \cdot \Delta t\right), \quad \lambda = 0.015\text{ hr}^{-1}
    \end{equation}
\end{itemize}

\begin{algorithm}[htbp]
\caption{FINDORA Multimodal Candidate Matching}
\label{alg:multimodal}
\begin{algorithmic}[1]
\REQUIRE Target report $I_t$, Candidate database $\mathcal{C}$, Weights $\mathbf{w}$
\ENSURE Ranked matches list $\mathcal{M}$
\STATE $\mathcal{M} \leftarrow \emptyset$
\STATE $\vec{v}_t \leftarrow \text{TF-IDF}(I_t.\text{title} + \text{`` ''} + I_t.\text{description})$
\STATE $\vec{c}_t \leftarrow \text{HexToRGB}(I_t.\text{color})$
\FOR{each candidate $I_j \in \mathcal{C}$}
    \IF{$I_t.\text{type} == I_j.\text{type}$}
        \STATE \textbf{continue} \COMMENT{Skip same polarity reports}
    \ENDIF
    \STATE $S_{\text{cat}} \leftarrow (I_t.\text{category} == I_j.\text{category}) \ ? \ 1.0 : 0.0$
    \STATE $\vec{v}_j \leftarrow \text{TF-IDF}(I_j.\text{title} + \text{`` ''} + I_j.\text{description})$
    \STATE $S_{\text{sem}} \leftarrow (\vec{v}_t \cdot \vec{v}_j) / (\|\vec{v}_t\|_2 \|\vec{v}_j\|_2)$
    \STATE $S_{\text{vis}} \leftarrow 1.0 - \|\vec{c}_t - \text{HexToRGB}(I_j.\text{color})\|_2 / 441.67$
    \STATE $\Delta d \leftarrow \text{Haversine}(I_t.\text{coords}, I_j.\text{coords})$
    \STATE $S_{\text{spa}} \leftarrow \exp(-\Delta d / 250.0)$
    \STATE $\Delta t \leftarrow |I_t.\text{time} - I_j.\text{time}|$ in hours
    \STATE $S_{\text{tem}} \leftarrow \exp(-0.015 \cdot \Delta t)$
    \STATE $\Phi \leftarrow w_v S_{\text{vis}} + w_s S_{\text{sem}} + w_c S_{\text{cat}} + w_l S_{\text{spa}} + w_t S_{\text{tem}}$
    \IF{$\Phi \ge 0.45$}
        \STATE $\mathcal{M} \leftarrow \mathcal{M} \cup \{(I_j, \Phi)\}$
    \ENDIF
\ENDFOR
\STATE Sort $\mathcal{M}$ in descending order by $\Phi$
\RETURN $\mathcal{M}$
\end{algorithmic}
\end{algorithm}

\section{Smart Verification and Handover Protocol}

\subsection{Zero-Knowledge Blind Challenge Verification}
To prevent fraudulent claims, private identifying traits (e.g., lock screen wallpaper, internal keychain color, phone case brand, serial digits) are kept hidden from public view. 

When a student clicks ``Claim Item'', the system generates interactive challenge questions:
\begin{enumerate}
    \item The claimant submits answers without knowing the true answers stored in the database.
    \item The server evaluates response similarity using token matching and Levenshtein edit distance:
    \begin{align}
        \text{Sim}(a_k, \kappa_k) &= 0.5 \left(\frac{|T(a_k) \cap T(\kappa_k)|}{|T(a_k) \cup T(\kappa_k)|}\right) \nonumber \\
        &\quad + 0.5 \left(1 - \frac{\text{Lev}(a_k, \kappa_k)}{\max(|a_k|, |\kappa_k|)}\right)
    \end{align}
    \item If the average score $\Omega \ge 0.70$, the claim is automatically approved; if $0.40 \le \Omega < 0.70$, it is sent for officer review; if $\Omega < 0.40$, it is rejected.
\end{enumerate}

\subsection{Behavioral Anti-Fraud Velocity Engine}
To stop individuals from repeatedly guessing answers, FINDORA AI monitors claimant behavior:
\begin{equation}
    \Psi = \min\left(100, \; 15 \times N_{\text{claims}}(24\text{h}) + 20 \times N_{\text{fails}}\right)
\end{equation}
If a claimant's risk score $\Psi \ge 70$, their account is temporarily restricted, and a security alert is triggered on the administrator dashboard.

\subsection{One-Time Handover Code Protocol}
Once a claim is verified, FINDORA AI creates a unique 6-character alphanumeric pickup code $K_{\text{code}} \in \{\text{A-Z}, 0-9\}^6$ and sends it directly to the verified owner's college email. When the student arrives at the campus security desk:
\begin{itemize}
    \item The student shows the email code to the officer.
    \item The officer enters the code into the FINDORA Admin Dashboard.
    \item If the code matches, the item status updates to \textbf{RECOVERED}, ensuring a reliable and recorded handover.
\end{itemize}

\section{Telegram Community Bot Integration}
To ensure all campus members are notified immediately when an item is found, FINDORA AI connects to a Telegram bot (\texttt{@findoravsb\_bot}) and channel (\texttt{findoravsbec}).
\begin{itemize}
    \item \textbf{Instant Broadcast:} When a new item is logged, an alert is pushed to the campus Telegram channel.
    \item \textbf{Bot Commands:} Students can type \texttt{/summary} to view current statistics, \texttt{/lost} to report lost items, and \texttt{/status [ID]} to check recovery progress.
\end{itemize}

\section{Experimental Results and Discussion}
FINDORA AI was tested in a 30-day pilot study at \textbf{V.S.B. Engineering College (VSBEC)} across major campus locations including the Central Library, Mechanical and Electrical Labs, Academic Blocks, and the Cafeteria. A dataset of 450 real and simulated lost-and-found items was evaluated.

\subsection{Matching Accuracy and Precision}
We evaluated FINDORA AI against traditional database keyword searches (SQL \texttt{LIKE}) and basic text-only models across key information retrieval metrics:
\begin{itemize}
    \item \textbf{Precision@$k$ (P@$k$):} Percentage of relevant matching items in the top $k$ recommendations.
    \item \textbf{Recall@5 (R@5):} Percentage of true matching items retrieved within the top 5 suggestions.
    \item \textbf{Mean Reciprocal Rank (MRR):} The average of the reciprocal ranks of the first correct match.
\end{itemize}

\begin{table}[htbp]
\caption{Comparison of Candidate Retrieval Accuracy}
\label{tab:retrieval_performance}
\centering
\setlength{\tabcolsep}{4pt}
\begin{tabular}{lcccc}
\toprule
\textbf{Model / Approach} & \textbf{P@1} & \textbf{P@3} & \textbf{R@5} & \textbf{MRR} \\
\midrule
Keyword Match (SQL \texttt{LIKE}) & 34.2\% & 48.6\% & 52.1\% & 0.412 \\
Text-Only TF-IDF Vector Space & 68.3\% & 79.1\% & 83.4\% & 0.741 \\
Semantic Embeddings Only & 74.5\% & 84.2\% & 87.9\% & 0.795 \\
\textbf{FINDORA AI (Multimodal Fusion)} & \textbf{88.6\%} & \textbf{94.2\%} & \textbf{97.8\%} & \textbf{0.914} \\
\bottomrule
\end{tabular}
\end{table}

As shown in Table~\ref{tab:retrieval_performance}, FINDORA AI achieves a \textbf{94.2\% Precision@3} and an \textbf{MRR of 0.914}, significantly outperforming standard keyword search (48.6\% P@3). Combining visual color metrics and spatial proximity helped resolve ambiguity between items with similar descriptions located in different departments.

\subsection{Operational Recovery Efficiency}
To evaluate real-world impact, we compared the recovery outcomes under FINDORA AI with historical manual logbook data.

\begin{table}[htbp]
\caption{Campus Operational Recovery Metrics}
\label{tab:operational_metrics}
\centering
\resizebox{\columnwidth}{!}{%
\setlength{\tabcolsep}{3pt}
\begin{tabular}{lccc}
\toprule
\textbf{Evaluation Metric} & \textbf{Manual Logbook} & \textbf{FINDORA AI} & \textbf{Improvement} \\
\midrule
Average Recovery Time & 72.4 hours & \textbf{4.2 hours} & \textbf{17.2$\times$ faster} \\
Median Recovery Time & 48.0 hours & \textbf{1.8 hours} & \textbf{26.6$\times$ faster} \\
Overall Campus Recovery Rate & 17.6\% & \textbf{84.3\%} & \textbf{+379\% increase} \\
False Claim Incidents & 14.2\% & \textbf{0.06\%} & \textbf{99.4\% drop} \\
Staff Workload (Hours/Week) & 28.5 hrs & \textbf{1.8 hrs} & \textbf{93.7\% saved} \\
Student Participation Rate & 4.1\% & \textbf{76.8\%} & \textbf{18.7$\times$ increase} \\
\bottomrule
\end{tabular}%
}
\end{table}

As indicated in Table~\ref{tab:operational_metrics}, FINDORA AI reduced average recovery time from \textbf{72.4 hours to 4.2 hours}. The campus recovery rate jumped from \textbf{17.6\% to 84.3\%}, while false claims dropped by 99.4\% due to the blind verification challenges.

\subsection{Zone-Wise Analysis and Student Feedback}
During the pilot, the highest volume of reported lost items occurred in the Central Library (34.2\%) and the Campus Cafeteria (27.8\%). The spatial decay formula prevented mismatching between items of the same color located in different campus blocks. Student surveys revealed that 91.4\% of users found the blind question verification intuitive and appreciated receiving instant notifications via Telegram.

\section{Conclusion and Future Work}
In this project, Team Techsquad developed \textbf{FINDORA AI} to solve the long-standing problem of lost item recovery in college environments. By integrating live camera GPS watermarking, multimodal matching (TF-IDF text similarity, normalized RGB color distance, and spatio-temporal decay), zero-knowledge blind challenge quizzes, and secure one-time handover codes, the system provides an end-to-end smart solution. The pilot testing at V.S.B. Engineering College demonstrated high retrieval accuracy (\textbf{94.2\% Precision@3}) and reduced return times by \textbf{17.2$\times$}.

In future work, we plan to implement edge-AI object recognition directly on mobile devices and explore automated smart locker integrations for 24/7 contactless item retrieval.

\section*{Acknowledgment}
The authors thank the faculty and administration of V.S.B. Engineering College (VSBEC) for their support and guidance throughout the development and pilot testing of this project.

\balance

\begin{thebibliography}{00}
\footnotesize
\bibitem{devlin2019bert} J. Devlin, M.-W. Chang, K. Lee, and K. Toutanova, ``BERT: Pre-training of Deep Bidirectional Transformers for Language Understanding,'' in \textit{Proc. NAACL-HLT}, 2019, pp. 4171--4186.
\bibitem{radford2021clip} A. Radford \textit{et al.}, ``Learning Transferable Visual Models From Natural Language Supervision (CLIP),'' in \textit{Proc. ICML}, 2021, pp. 8748--8763.
\bibitem{robertson2009bm25} S. Robertson and H. Zaragoza, ``The Probabilistic Relevance Framework: BM25 and Beyond,'' \textit{Found. Trends Inf. Retr.}, vol. 3, no. 4, pp. 333--389, 2009.
\bibitem{goldwasser1989zkp} S. Goldwasser, S. Micali, and C. Rackoff, ``The Knowledge Complexity of Interactive Proof Systems,'' \textit{SIAM J. Comput.}, vol. 18, no. 1, pp. 186--208, 1989.
\bibitem{levenshtein1966binary} V. I. Levenshtein, ``Binary codes capable of correcting deletions, insertions, and reversals,'' \textit{Soviet Physics Doklady}, vol. 10, no. 8, pp. 707--710, 1966.
\bibitem{deepmind2025gemini} Google DeepMind, ``Gemini 2.5: Multimodal Foundations for Next-Generation Scalable Reasoning,'' \textit{Google Research Tech. Rep.}, 2025.
\bibitem{vsbec2026problem} VSBEC Smart Campus Hackathon Committee, ``Problem Statement 22: Smart Lost and Found Management System Specifications,'' \textit{V.S.B. Engineering College}, Karur, Tech. Rep. PS-22, 2026.
\bibitem{manning2008ir} C. D. Manning, P. Raghavan, and H. Sch\"{u}tze, \textit{Introduction to Information Retrieval}. Cambridge: Cambridge Univ. Press, 2008.
\bibitem{haversine1984sinot} R. W. Sinnott, ``Virtues of the Haversine,'' \textit{Sky and Telescope}, vol. 68, no. 2, p. 159, 1984.
\end{thebibliography}

\end{document}
